
\section{Infinite $(8,7)$ families}\label{sec:families}

This section proves the four parts of Theorem~\ref{main:families}. Three
mechanisms are involved: a diagonal retraction of a catalogued six-square
carrier family, a bridge completion inside a catalogued carrier type, and
an elliptic-orbit lift of one fixed carrier.

\subsection{The retraction family and its least-shift law}

The carrier is the six-square fully magic family recorded on Boyer's
catalogue pages~\cite{BOYWEB} and attributed there to Wesolowski; the
catalogue states explicitly that the recurrence below ``was not given by
Arkadiusz'', and the $n=1$ grid already occurs as Configuration~II of
Bremner~\cite{BRE01}, so we speak of the Bremner/Wesolowski carrier. With
$(x_0,y_0,z_0)=(3,1,1)$ and
\[
x_{n+1}=3x_n+2z_n,\qquad y_{n+1}=x_n,\qquad z_{n+1}=x_n+z_n,
\]
the carrier $W_n$ is fully magic with centre $E=4z^4+5z^2+1$ and common sum
$3E$, where $(x,y,z)=(x_n,y_n,z_n)$. Subtracting $t=6z^2$ from its
antidiagonal and using the recurrence identities
\[
x^2+y^2z^2-6z^2=(yz+1)^2,\qquad
x^2z^2+y^2-6z^2=(xz-1)^2,\qquad
\tfrac{(x^2+y^2)(z^2+1)}{2}-6z^2=4z^4-z^2+1
\]
produces
\[
N_n=\begin{pmatrix}
(xy-z)^2 & (xz+y)^2 & (yz+1)^2\\
(x+yz)^2 & 4z^4-z^2+1 & (xz-y)^2\\
(xz-1)^2 & (yz-x)^2 & (xy+z)^2
\end{pmatrix}.
\]

\begin{proof}[Proof of Theorem~\ref{main:families}(i)]
By Lemma~\ref{lem:retraction}(i) the seven unselected lines of $N_n$ share
the sum $3E-6z^2=3(4z^4+3z^2+1)$, the antidiagonal sums to
$3E-18z^2=3(4z^4-z^2+1)$, and the defect is $12z^2$. For $z>1$,
\[
(2z^2-1)^2<4z^4-z^2+1<(2z^2)^2,
\]
so the centre is not a square. The eight positive roots are, in increasing
order,
\[
yz-x<yz+1<x+yz<xy-z<xy+z<xz-y<xz-1<xz+y:
\]
positivity of the first follows from $yz-x=y(z-1)-2z\ge3(z-1)-2z=z-3>0$,
most gaps are immediate, and the two non-obvious ones reduce, through the
recurrence invariants, to
\[
(xy-z)-(x+yz)=2z^2-yz-y-3z+1>0,\qquad
(xz-y)-(xy+z)=(y-1)(z-1)-2>0,
\]
the first bounded below by $(z-1)(z-2)>0$ using $y\le z-1$, the second at
least $4$. Hence all nine entries are positive and distinct. For
root-primitivity: an integer dividing all eight roots divides both
$(yz+1)-(yz-x)=x+1$ and $(x+yz)-(yz+1)=x-1$, hence divides $2$; the
recurrence keeps $x,y$ odd and toggles the parity of $z$, so for odd $n$
all eight roots are odd and their gcd is $1$. There are infinitely many
odd indices.
\end{proof}

\begin{proposition}[Complete pair-closing classification]\label{prop:factorpairs}
Fix $n$ and write the retracted antidiagonal entries of the carrier as
$C<E<G$ with common gap $\Delta$. Every positive shift $0<t<C$ for which at
least two of $C-t$, $E-t$, $G-t$ are squares arises from a same-parity
factor pair, and every retained pair works: for branch $CE$, from
$de=\Delta$, $0<d<e$, $d\equiv e\pmod2$, via $u=(e-d)/2$, $w=(e+d)/2$,
$t=C-u^2$; for branch $EG$, from $de=\Delta$ via $w=(e-d)/2$,
$v=(e+d)/2$, $t=E-w^2$; for branch $CG$ from $de=2\Delta$
via $u=(e-d)/2$, $v=(e+d)/2$, $t=C-u^2$.
\end{proposition}

\begin{proof}
For $CE$, subtracting the two square equations gives
$w^2-u^2=\Delta=(w-u)(w+u)$, so $d=w-u$ and $e=w+u$ have the stated product
and parity, and the inverse formulas recover $u,w$. The other branches are
identical with difference $\Delta$ or $2\Delta$. This is a divisor
classification, not a scan over $t$.
\end{proof}

\begin{proposition}[Least shift]\label{prop:leastshift}
For every $n\ge2$, $t_0=6z^2$ is the least positive shift closing at least
one pair; at $n=1$ the least is $9$. Hence $t_{\min}(1)=9$ and
$t_{\min}(n)=6z_n^2$ for $n\ge2$. Leastness does not mean uniqueness:
larger nondegenerate closing shifts exist and are recorded exactly.
\end{proposition}

\begin{proof}[Proof frame]
For $n\ge2$ the recurrence on $(y,z)$ is $(y,z)\mapsto(y+2z,y+3z)$,
starting at $(11,15)$, and direct induction gives
\[
z>y,\quad 2y>z,\quad 5y>3z,\quad 5z>6y,\quad 8y>5z,\quad z\ge15 .
\]
Assume $0<t<t_0$ closes a pair and put $h=t_0-t$, $0<h<6z^2$. At $t_0$ the
three entries are $C-t_0=Y^2$ with $Y=yz+1$, $E-t_0=4z^4-z^2+1$, and
$G-t_0=X^2$ with $X=xz-1$; each entry can be square inside the short
interval only along an explicit one-parameter list (for the $C$ entry,
$C-t=(Y+a)^2$ forces $h=a(2Y+a)$, and similarly for $E$ and $G$), and the
listed windows are excluded pairwise by the induction inequalities. The
complete three-case interval analysis, with its exact certificate, is
reproduced in Appendix~B. At $n=1$, $(C,E,G)=(265,1105,1945)$: for
$1\le t\le8$ each retracted entry lies strictly between consecutive squares
($16^2<C-t<17^2$, $33^2<E-t<34^2$, $44^2<G-t<45^2$), so no pair closes
below $t=9$, while $t=9$ gives $(256,1096,1936)=(16^2,1096,44^2)$.
\end{proof}

Parts (i)--(ii) of Theorem~\ref{main:families} are proved.

\subsection{The bridge completion in type 6.VI}

The catalogued carrier type \texttt{6:6} (6.VI here) is normalized, after
the dihedral action, as
\[
M(A,N,Z)=\begin{pmatrix}
2N-Z & A & N-A+Z\\
2Z-A & N & 2N-2Z+A\\
N+A-Z & 2N-A & Z
\end{pmatrix},
\]
with square cells $A=a^2$, $N=n^2$, $Z=z^2$, $B=2N-A=b^2$, $D=2Z-A=d^2$,
$C=2N-D=c^2$. The type taxonomy and database records are Fituvalu's~\cite{FIT}; the
fixed numerical carrier used below is historically Bremner's
Configuration~VI~\cite{BRE01}.

\begin{proposition}[The four-progression network]\label{prop:fourap}
The six source squares satisfy $a^2+b^2=2n^2$, $d^2+c^2=2n^2$ and
$d^2+a^2=2z^2$, so the source contains the three square progressions
$(a^2,n^2,b^2)$, $(d^2,n^2,c^2)$, $(d^2,z^2,a^2)$. With
$\delta=a^2-z^2=z^2-d^2>0$, the selected antidiagonal is identically
$(n^2-\delta,\,n^2,\,n^2+\delta)$, and subtracting $t>0$ closes it to
squares $(x^2,y^2,w^2)$ exactly when
\[
x^2=y^2-\delta,\qquad w^2=y^2+\delta,\qquad t=n^2-y^2>0:
\]
a \emph{fourth} square progression, with the difference of the bridge
$(d^2,z^2,a^2)$.
\end{proposition}

\begin{proof}
Substituting $D=2Z-A$ into the normalized array gives $A-Z=Z-D=\delta$, and
the antidiagonal entries become $N-\delta$, $N$, $N+\delta$; a common
subtraction preserves their differences, and the closure conditions are the
displayed ones.
\end{proof}

\begin{proposition}[Complete bridge-completion classification]\label{prop:bridge}
Fix an increasing bridge progression $(d^2,z^2,a^2)$ of difference
$\delta$. Every positive type-6.VI completion arises from a same-parity
factor pair
\[
uv=2\delta,\qquad 0<u<v,\qquad u\equiv v\pmod 2,
\]
via $b=(v-u)/2$, $c=(v+u)/2$, $n^2=(a^2+b^2)/2$, a pair being retained
exactly when $n^2$ is an integer square; then $d^2+c^2=a^2+b^2=2n^2$
automatically. The pair $(u,v)=(a-d,a+d)$ always returns the tautological
completion $(b,n,c)=(d,z,a)$, which repeats square cells and is not
admissible. Positivity, nine-way distinctness and the six-square mask are
separate final checks.
\end{proposition}

\begin{proof}
Subtracting the two centre relations gives $c^2-b^2=a^2-d^2=2\delta$, so
$u=c-b$, $v=c+b$ have the stated product and parity, and inversion recovers
$b,c$; conversely any retained pair satisfies $c^2-b^2=2\delta=a^2-d^2$, so
the two expressions for $2n^2$ agree.
\end{proof}

\begin{proof}[Proof of Theorem~\ref{main:families}(iii)]
For the fixed bridge $(23^2,37^2,47^2)$, $\delta=840$: running
Proposition~\ref{prop:bridge} over the finitely many same-parity factor
pairs of $2\delta=1680$ retains exactly the tautological $(23,37,47)$ and
the unique admissible completion $(79,65,89)$ (exact finite check;
certificate in Appendix~B). Coupling the resulting carrier with Robertson's
classical progression $(1^2,29^2,41^2)$ of the same
difference~\cite{ROB96} retracts the antidiagonal by $\tau=65^2-29^2=3384$
and yields the explicit $(8,7)$ array
\[
\begin{pmatrix}
7081 & 2209 & 1\\
529 & 841 & 7921\\
1681 & 6241 & 1369
\end{pmatrix}
\]
with seven-line sum $9291$, failed antidiagonal sum $2523$, and unique
nonsquare $7081$.
\end{proof}

\subsection{The elliptic-orbit lift}

The fixed carrier is
\[
M_0=\begin{pmatrix}
7081 & 2209 & 3385\\
529 & 4225 & 7921\\
5065 & 6241 & 1369
\end{pmatrix},
\]
Bremner's Configuration~VI~\cite{BRE01}, in the normalized orientation of
Fituvalu database row 340545~\cite{FIT}, is fully magic with sum $12675$.
Its six square entries have outside roots
\[
\{23,37,47,79,89\},
\]
and its selected antidiagonal is
$(65^2-840,\,65^2,\,65^2+840)$. Let
\[
E\colon\;Y^2=X^3-840^2X,\qquad Q=(1960,78400),\qquad P=2Q=(841,-1189).
\]

\begin{lemma}[Doubling produces a square progression]\label{lem:doubling}
For $R=(u,v)\in E(\Q)$ with $v\neq0$,
\[
x(2R)=\Bigl(\frac{u^2+840^2}{2v}\Bigr)^{2},\qquad
x(2R)\mp840=\Bigl(\frac{u^2\mp1680u-840^2}{2v}\Bigr)^{2},
\]
so every finite point of $2E(\Q)$ supplies three rational squares in
arithmetic progression with common difference $840$ --- the classical
congruent-number correspondence (cf.~\cite{ROB96}), in the exact convention
used by the lift.
\end{lemma}

\begin{lemma}[Fixed-carrier lift]\label{lem:lift}
Let $R\in2E(\Q)$ with $840<x(R)<65^2$. Write the positive roots of
$x(R)-840$, $x(R)$, $x(R)+840$ in lowest terms and let $\lambda$ be the
\emph{least} common multiple of their denominators. Scale $M_0$ by
$\lambda^2$ and subtract $\tau=\lambda^2(65^2-x(R))$ from the selected
antidiagonal; the antidiagonal becomes the three cleared squares. If the
three new roots avoid $\{23,37,47,79,89\}$, the output is positive with
nine distinct entries and exactly eight squares, the unique nonsquare being
$7081\lambda^2$ --- and it is \emph{entry-primitive}: a common prime $p$ of
all entries would divide $\lambda$; choosing a reduced root denominator of
maximal $p$-adic valuation, its cleared numerator stays prime to $p$, and
its square blocks $p$. The argument uses minimality of $\lambda$ and fails
for an arbitrary common multiple.
\end{lemma}

\begin{lemma}[$Q$ is non-torsion]\label{lem:nontorsion}
$P=2Q=(841,-1189)$ is integral with $1189=29\cdot41$, while
$4\cdot840^{6}$ has only the prime divisors $2,3,5,7$; by the Lutz--Nagell
criterion~\cite[Chap.~VII, Prop.~3.1]{SIL09} an integral torsion point with nonzero ordinate would satisfy
$Y^2\mid4\cdot840^6$, so $P$ is non-torsion, hence so is $Q$.
Independently, $\#E(\mathbb F_{11})=12$ and $\#E(\mathbb F_{13})=20$ at
good primes bound the rational torsion order by $\gcd(12,20)=4$, and $P$
is neither infinity nor rational two-torsion.
\end{lemma}

\begin{proof}[Proof of Theorem~\ref{main:families}(iv)]
$P$ lies in the identity component $E(\mathbb R)^{0}$, a circle group, and
is non-torsion, so its positive multiples are dense there. The open set
$U=\{R\in E(\mathbb R)^{0}:840<x(R)<4225\}$ contains $P$; removing the
finite set where one of the three roots equals one of $23,37,47,79,89$
leaves a nonempty open set $V$ --- the roots at $P$ itself are $(1,29,41)$.
Density gives infinitely many positive multiples $mP=2mQ$ in $V$, and
Lemma~\ref{lem:lift}, always with the least common multiple, turns each
return into a primitive positive $(8,7)$ array. If the outputs at indices
$m,n$ were related by a $D_4$ action and a positive square scale $\rho^2$,
the unique nonsquare entries would force
$7081\lambda_n^2=\rho^2\,7081\lambda_m^2$, so $\rho=\lambda_n/\lambda_m$;
a $D_4$ action fixes the centre, so comparing centres gives
$x(mP)=x(nP)$, hence $mP=\pm nP$, and non-torsion with $m,n>0$ forces
$m=n$. Distinct return indices therefore give distinct
square-scale/$D_4$ classes. Two exact witnesses, including the $m=3$
return, are frozen in the companion repository.
\end{proof}

\begin{corollary}[Positive-density orbit lift]\label{cor:orbitdensity}
Let $A(N)$ count the indices $1\le m\le N$ for which $mP$ lies in the
admissible set $V$ of the preceding proof. Then
\[
A(N)=\delta N+o(N),
\qquad
\delta=
\frac{\displaystyle\int_{840}^{4225}
dx/\sqrt{x^3-840^2x}}
{\displaystyle\int_{840}^{\infty}dx/\sqrt{x^3-840^2x}}
=1-I_{(840/4225)^2}\!\left(\frac14,\frac12\right),
\]
and
\[
\delta=0.6585271498271213\ldots .
\]
The same asymptotic counts the square-scale/$D_4$-inequivalent primitive
$(8,7)$ classes produced among $P,2P,\ldots,NP$.
\end{corollary}

\begin{proof}
The multiples of the non-torsion point $P$ are equidistributed in the
circle group $E(\mathbb R)^0$ with respect to normalized Haar measure; see
Cowan~\cite{COW20} for the natural-density formulation for real elliptic
curves. The boundary of $840<x<4225$ has measure zero, and deleting the
finitely many collision points changes neither measure nor density. On the
identity component the invariant measure is proportional to
$dx/\sqrt{x^3-840^2x}$, giving the ratio above. Substitution
$x=840/t^2$ changes it to
\[
\frac{\displaystyle\int_{\sqrt{840/4225}}^1dt/\sqrt{1-t^4}}
{\displaystyle\int_0^1dt/\sqrt{1-t^4}},
\]
which is the stated regularized incomplete-beta value. The final counting
claim follows from the inequivalence argument in the preceding proof.
\end{proof}

\begin{remark}[The universal lemniscatic law behind $\delta$]
The special-function value in Corollary~\ref{cor:orbitdensity} is not
peculiar to $840$. For $d>0$ let
\[
E_d:\quad Y^2=X^3-d^2X.
\]
If $R=(X,Y)$ is distributed according to normalized Haar probability on
$E_d(\mathbb R)^0$ and $U=(d/X)^2$, then the push-forward density is
\[
 d\mu_U(u)=
 \frac{u^{-3/4}(1-u)^{-1/2}}
 {\mathrm B(\tfrac14,\tfrac12)}\,du,
 \qquad 0<u<1;
\]
equivalently, $U\sim\mathrm{Beta}(\tfrac14,\tfrac12)$. Indeed,
$X=du^{-1/2}$ sends the invariant differential
$dX/\sqrt{X^3-d^2X}$ to a constant multiple of
$u^{-3/4}(1-u)^{-1/2}du$. Consequently, for every $B>d$,
\[
 \mu\{R\in E_d(\mathbb R)^0:d<X(R)<B\}
 =1-I_{(d/B)^2}\!\left(\frac14,\frac12\right).
\]
The normalizing integral is the classical lemniscatic value
\[
 \int_0^1\frac{dt}{\sqrt{1-t^4}}
 =\frac{\Gamma(1/4)^2}{4\sqrt{2\pi}}
\]
recorded in the NIST Digital Library of Mathematical
Functions~\cite[\S19.20(i), Eq.~(19.20.2)]{DLMF}. This universality reflects
$j(E_d)=1728$: after real scaling, every curve in the family has the same
normalized invariant measure, so the probability depends only on $B/d$.
For a non-torsion point in $E_d(\mathbb R)^0$, equidistribution transfers
this probability to orbit indices. A point on the other real component
requires separate bookkeeping for its even and odd multiples. Thus the
remark explains the density constant but adds no new arithmetic input to
the corollary.
\end{remark}

\begin{remark}[From constructions to the wall]
Part (iv) reaches the $(8,7)$ metric infinitely often from one carrier.
Closing the eighth line of the same transversal frame is the interaction
problem of Section~\ref{sec:surface}: by Remark~\ref{rem:fourthAP} it
demands a fourth square progression through the centres,
and that single extra condition changes the landscape from constructive to
arithmetic.
\end{remark}
